%================================ REFERENCES ==================================
\clearpage
\addcontentsline{toc}{chapter}{References}
\renewcommand{\bibname}{References}

\begin{thebibliography}{00}
\addtolength{\itemsep}{-2pt}

\bibitem{bloom1984} B.~S. Bloom, ``The 2 sigma problem: The search for methods
of group instruction as effective as one-to-one tutoring,'' \emph{Educational
Researcher}, vol.~13, no.~6, pp.~4--16, 1984.

\bibitem{anderson1995cognitive} J.~R. Anderson, A.~T. Corbett, K.~R. Koedinger,
and R.~Pelletier, ``Cognitive tutors: Lessons learned,'' \emph{The Journal of
the Learning Sciences}, vol.~4, no.~2, pp.~167--207, 1995.

\bibitem{vanlehn2011} K.~VanLehn, ``The relative effectiveness of human
tutoring, intelligent tutoring systems, and other tutoring systems,''
\emph{Educational Psychologist}, vol.~46, no.~4, pp.~197--221, 2011.

\bibitem{kasneci2023chatgpt} E.~Kasneci \emph{et al.}, ``ChatGPT for good? On
opportunities and challenges of large language models for education,''
\emph{Learning and Individual Differences}, vol.~103, 102274, 2023.

\bibitem{paivio1991} A.~Paivio, ``Dual coding theory: Retrospect and current
status,'' \emph{Canadian Journal of Psychology}, vol.~45, no.~3, pp.~255--287,
1991.

\bibitem{sweller1988} J.~Sweller, ``Cognitive load during problem solving:
Effects on learning,'' \emph{Cognitive Science}, vol.~12, no.~2, pp.~257--285,
1988.

\bibitem{mayer2014multimedia} R.~E. Mayer, Ed., \emph{The Cambridge Handbook of
Multimedia Learning}, 2nd~ed. Cambridge, U.K.: Cambridge Univ. Press, 2014.

\bibitem{mayermoreno2003} R.~E. Mayer and R.~Moreno, ``Nine ways to reduce
cognitive load in multimedia learning,'' \emph{Educational Psychologist},
vol.~38, no.~1, pp.~43--52, 2003.

\bibitem{vaswani2017attention} A.~Vaswani \emph{et al.}, ``Attention is all you
need,'' in \emph{Proc. Advances in Neural Information Processing Systems
(NeurIPS)}, 2017, pp.~5998--6008.

\bibitem{brown2020gpt3} T.~B. Brown \emph{et al.}, ``Language models are
few-shot learners,'' in \emph{Proc. Advances in Neural Information Processing
Systems (NeurIPS)}, 2020, pp.~1877--1901.

\bibitem{gemini2023} Gemini Team, Google, ``Gemini: A family of highly capable
multimodal models,'' arXiv:2312.11805, 2023.

\bibitem{chen2021codex} M.~Chen \emph{et al.}, ``Evaluating large language
models trained on code,'' arXiv:2107.03374, 2021.

\bibitem{liu2023promptsurvey} P.~Liu \emph{et al.}, ``Pre-train, prompt, and
predict: A systematic survey of prompting methods in natural language
processing,'' \emph{ACM Computing Surveys}, vol.~55, no.~9, pp.~1--35, 2023.

\bibitem{ji2023hallucination} Z.~Ji \emph{et al.}, ``Survey of hallucination in
natural language generation,'' \emph{ACM Computing Surveys}, vol.~55, no.~12,
pp.~1--38, 2023.

\bibitem{rombach2022stable} R.~Rombach, A.~Blattmann, D.~Lorenz, P.~Esser, and
B.~Ommer, ``High-resolution image synthesis with latent diffusion models,'' in
\emph{Proc. IEEE/CVF Conf. Computer Vision and Pattern Recognition (CVPR)},
2022, pp.~10684--10695.

\bibitem{wei2022cot} J.~Wei \emph{et al.}, ``Chain-of-thought prompting elicits
reasoning in large language models,'' in \emph{Proc. Advances in Neural
Information Processing Systems (NeurIPS)}, 2022.

\bibitem{yao2023react} S.~Yao \emph{et al.}, ``ReAct: Synergizing reasoning and
acting in language models,'' in \emph{Proc. Int. Conf. Learning Representations
(ICLR)}, 2023.

\bibitem{madaan2023selfrefine} A.~Madaan \emph{et al.}, ``Self-Refine: Iterative
refinement with self-feedback,'' in \emph{Proc. Advances in Neural Information
Processing Systems (NeurIPS)}, 2023.

\bibitem{shinn2023reflexion} N.~Shinn \emph{et al.}, ``Reflexion: Language
agents with verbal reinforcement learning,'' in \emph{Proc. Advances in Neural
Information Processing Systems (NeurIPS)}, 2023.

\bibitem{chen2023selfdebug} X.~Chen, M.~Lin, N.~Sch\"arli, and D.~Zhou,
``Teaching large language models to self-debug,'' arXiv:2304.05128, 2023.

\bibitem{langchain} H.~Chase, ``LangChain,'' open-source software, 2022.
[Online]. Available: \url{https://github.com/langchain-ai/langchain}

\bibitem{langgraph} LangChain, ``LangGraph: Building stateful, multi-actor
applications with LLMs,'' software documentation. [Online]. Available:
\url{https://langchain-ai.github.io/langgraph/}

\bibitem{svgspec} World Wide Web Consortium (W3C), ``Scalable Vector Graphics
(SVG) 1.1 (Second Edition),'' W3C Recommendation, 2011. [Online]. Available:
\url{https://www.w3.org/TR/SVG11/}

\bibitem{carlier2020deepsvg} A.~Carlier, M.~Danelljan, A.~Alahi, and
R.~Timofte, ``DeepSVG: A hierarchical generative network for vector graphics
animation,'' in \emph{Proc. Advances in Neural Information Processing Systems
(NeurIPS)}, 2020.

\bibitem{rodriguez2023starvector} J.~A. Rodriguez \emph{et al.}, ``StarVector:
Generating scalable vector graphics code from images,'' arXiv:2312.11556, 2023.

\bibitem{oord2016wavenet} A.~van den Oord \emph{et al.}, ``WaveNet: A generative
model for raw audio,'' arXiv:1609.03499, 2016.

\bibitem{shen2018tacotron2} J.~Shen \emph{et al.}, ``Natural TTS synthesis by
conditioning WaveNet on mel spectrogram predictions,'' in \emph{Proc. IEEE Int.
Conf. Acoustics, Speech and Signal Processing (ICASSP)}, 2018, pp.~4779--4783.

\bibitem{offlinefirst} A.~Feyerke, ``Designing offline-first web apps,'' \emph{A
List Apart}, no.~373, 2013. [Online]. Available:
\url{https://alistapart.com/article/offline-first/}

\bibitem{androidguide} Google, ``Guide to app architecture,'' Android Developers
documentation. [Online]. Available:
\url{https://developer.android.com/topic/architecture}

\bibitem{martin2017clean} R.~C. Martin, \emph{Clean Architecture: A Craftsman's
Guide to Software Structure and Design}. Boston, MA: Prentice Hall, 2017.

\bibitem{twelvefactor} A.~Wiggins, ``The Twelve-Factor App,'' 2017. [Online].
Available: \url{https://12factor.net/}

\bibitem{gamma1994} E.~Gamma, R.~Helm, R.~Johnson, and J.~Vlissides,
\emph{Design Patterns: Elements of Reusable Object-Oriented Software}. Reading,
MA: Addison-Wesley, 1994.

\end{thebibliography}
